\begin{frame}[plain]
\frametitle{Retrieval rises sharply with response length}
\begin{tikzpicture}[remember picture,overlay]
  \centerfigure{1.55}{\begin{tikzpicture}[text width=,anchor=center]
\path[use as bounding box] (0,0) rectangle (14.6,5.8);
\begin{axis}[
  tallslideaxis,
  xmin=-0.15,xmax=3.15,ymin=0,ymax=110,
  xtick={0,1,2,3},xticklabels={1--2,3--5,6--20,21+},
  ytick={0,20,40,60,80,100},
  legend style={anchor=north west,at={(0.025,0.98)},draw=none,fill=white,
    font=\fs{9}{12},cells={anchor=west}},
]
\addplot[SlideBlue,line width=1.2pt,mark=*,mark size=2.2pt,
  error bars/.cd,y dir=both,y explicit,
  error bar style={line width=0.8pt},
  error mark options={rotate=90,mark size=3pt}]
  table[x=x,y=en_pt,y error minus=en_pt_minus,y error plus=en_pt_plus,
    col sep=tab]{../shared/data/length-results.tsv};
\addlegendentry{English query $\rightarrow$ Portuguese candidates}
\addplot[MidGray,dashed,line width=1.2pt,mark=square*,mark size=2.0pt,
  error bars/.cd,y dir=both,y explicit,
  error bar style={line width=0.8pt},
  error mark options={rotate=90,mark size=3pt}]
  table[x=x,y=pt_en,y error minus=pt_en_minus,y error plus=pt_en_plus,
    col sep=tab]{../shared/data/length-results.tsv};
\addlegendentry{Portuguese query $\rightarrow$ English candidates}
\draw[SlideBlue,line width=0.6pt]
  (axis cs:0.04,8.385707) -- (axis cs:0.35,8.385707);
\node[anchor=west,font=\fs{10}{12}\bfseries,text=SlideBlue]
  at (axis cs:0.40,8.385707) {8.4\%};
\node[anchor=south west,font=\fs{10}{12}\bfseries,text=MidGray] at (axis cs:0.08,24) {13.0\%};
\node[anchor=east,font=\fs{10}{12}\bfseries,text=SlideBlue] at (axis cs:2.96,106) {99.1\%};
\node[anchor=east,font=\fs{10}{12}\bfseries,text=MidGray] at (axis cs:2.96,91) {98.5\%};
\end{axis}
\plotlabels{Rank-one recovery (\%)}{Portuguese source length in words}
\end{tikzpicture}
}
  \takeaway{Short responses offer fewer cues for recovering their paired text.}
  \source{Camera-ready article, Figure 1. 95\% cluster-bootstrap intervals, 2,000 resamples over survey-specific records.}
\end{tikzpicture}
\note{
  - 04:45--06:15, 90 seconds\\
  - Orient the audience. The horizontal axis uses Portuguese source length, with four prespecified bands. The vertical axis shows correct rank-one recovery.\\
  - Blue is an English query searching Portuguese candidates. Dashed gray is a Portuguese query searching English candidates.\\
  - One- or two-word responses have recovery rates of about 8 to 13 percent. Responses longer than twenty words reach about 98 to 99 percent.\\
  - Error bars are 95 percent cluster-bootstrap intervals from 2,000 resamples. All length groups use the same full candidate corpus.\\
  - This is an association between length and recoverability in this corpus. Bilingual judgments would connect individual misses to translation quality.\\
  - Source: camera-ready article, Figure 1.
}
\end{frame}

\begin{frame}[plain]
\frametitle{Quality estimates follow a different pattern}
\begin{tikzpicture}[remember picture,overlay]
  \centerfigure{1.55}{\begin{tikzpicture}[text width=,anchor=center]
\path[use as bounding box] (0,0) rectangle (14.6,5.8);
\begin{axis}[
  tallslideaxis,
  xmin=-0.15,xmax=3.15,ymin=0,ymax=1,
  xtick={0,1,2,3},xticklabels={1--2,3--5,6--20,21+},
  ytick={0,0.25,0.5,0.75,1},
]
\addplot[SlideBlue,line width=1.2pt,mark=*,mark size=2.6pt,
  nodes near coords,
  every node near coord/.append style={font=\fs{11}{14}\bfseries,
    text=Ink,anchor=south,yshift=3pt,
    /pgf/number format/fixed,/pgf/number format/precision=3,
    /pgf/number format/fixed zerofill}]
  table[x=x,y=cometkiwi,col sep=tab]{../shared/data/length-results.tsv};
\end{axis}
\plotlabels{Mean COMETKiwi score}{Portuguese source length in words}
\end{tikzpicture}
}
  \takeaway{Bilingual human ratings would calibrate these quality estimates.}
  \source{Camera-ready article, length-stratified COMETKiwi means. COMETKiwi DA-XL, Rei et al., WMT 2023.}
\end{tikzpicture}
\note{
  - 06:15--07:25, 70 seconds\\
  - This plot keeps the same length bands and shows mean COMETKiwi scores on their model scale.\\
  - The mean is highest for the three-to-five-word group, then falls for the longer groups. Retrieval continued to rise across those groups.\\
  - The disagreement is informative. Retrieval measures distinguishability among candidates, while COMETKiwi estimates quality using learned source--translation relations.\\
  - The study has automated criteria rather than bilingual human ratings. Those ratings would be the next step for calibration in this setting.\\
  - Source: camera-ready article, Retrieval Varies Sharply with Response Length and Discussion. Exact means are 0.6414, 0.7798, 0.7526 and 0.7152. Display labels round to three decimals.
}
\end{frame}

\begin{frame}[plain]
\frametitle{Emotion-lexicon coverage}
\begin{tikzpicture}[remember picture,overlay]
  \sput{1.3}{1.62}{13.4}{\centering\fs{12}{15}
    About one quarter of responses have \textbf{zero covered emotion terms}.}
  \centerfigure{2.15}{\begin{tikzpicture}[text width=,anchor=center]
\path[use as bounding box] (0,0) rectangle (14.6,5.0);
\begin{axis}[
  slideaxis,scale only axis,width=12.1cm,height=3.0cm,
  at={(1.9cm,1.0cm)},anchor=south west,xbar,bar width=0.8cm,
  xmin=0,xmax=100,ymin=-0.6,ymax=1.6,y dir=reverse,
  ytick={0,1},yticklabels={Portuguese,English},
  xtick={0,25,50,75,100},
  ymajorgrids=false,xmajorgrids=true,
  legend style={at={(0.5,1.10)},anchor=south,legend columns=2,draw=none,
    /tikz/every even column/.append style={column sep=12pt},
    font=\fs{9}{12},cells={anchor=west}},
  bar shift=0pt,
]
\addplot[fill=Grid!60,draw=none,forget plot] coordinates {(100,0) (100,1)};
\addplot[fill=SlideBlue,draw=none,forget plot,
  nodes near coords={\pgfmathprintnumber[fixed,precision=2,fixed zerofill]{\pgfplotspointmeta}\%},
  every node near coord/.append style={anchor=east,xshift=-4pt,
    font=\fs{11}{14}\bfseries,text=white,
    /pgf/number format/fixed,/pgf/number format/precision=2,
    /pgf/number format/fixed zerofill}]
  table[x=zero_coverage,y expr=\coordindex,
    col sep=tab]{../shared/data/emotion-coverage.tsv};
\addlegendimage{area legend,fill=SlideBlue,draw=none}
\addlegendentry{Zero coverage}
\addlegendimage{area legend,fill=Grid!60,draw=none}
\addlegendentry{At least one covered emotion term}
\end{axis}
\node[plainlabel,font=\fs{10}{12}] at (7.95,0.18) {Responses (\%)};
\end{tikzpicture}
}
  \takeaway{Affective comparisons need explicit reporting of lexical coverage.}
  \source{Camera-ready article, Language and Affective Coverage. EmoAtlas, Semeraro et al., Behavior Research Methods, 2025.}
\end{tikzpicture}
\note{
  - 07:25--08:25, 60 seconds\\
  - The blue portion shows responses for which the language-specific lexicon returns zero emotion coverage: 24.67 percent in Portuguese and 23.48 percent in English.\\
  - Zero coverage describes the resource's output. A response can carry meaning or emotion beyond the vocabulary recognized by the lexicon.\\
  - Empty affective profiles affect distance calculations mechanically. The paper therefore also reports comparisons conditional on nonzero profiles in both languages.\\
  - This is a resource-coverage constraint to carry into interpretation, especially when translated text supports mental-health research.\\
  - Source: camera-ready article, Language and Affective Coverage. Gray bar segments are the complements of the reported zero-coverage rates.
}
\end{frame}
